\begin{frame}{}
    \centering\small
    \tikzset{block/.style={rectangle, draw=gray, fill=gray!10, thick, minimum width=3em, minimum height=2em, align=center}}
    \begin{tikzpicture}
        \node[block,] (BM) {重みなし\\二部マッチング};
        \node[block,above=1em of BM.north west, xshift=-2em] (WBM) {重みつき\\二部マッチング};
        \node[block,above=5em of BM.north east, xshift=1em] (NF) {最大流};
        \node[block,above=9em of BM] (MCF) {最小費用流};
        % \node[block,above=1em of NF.north east, xshift=1em] (MC) {最小カット};
        \node[block,right=11em of BM] (MTR) {マトロイド};
        \node[block,above=7em of MTR] (MI) {マトロイド交差};
        \node[block,above=3em of MCF, xshift=7em] (SFM) {劣モジュラ最小化};
        \graph[mygraph]{
            (BM) -> {(WBM), (NF)} -> (MCF);
            {(MTR), (BM)} -> (MI);
            % (NF) -> (MC);
            {(MCF), (MI)} -> (SFM);
        };
        \begin{scope}[on background layer]
            \node[mysubset, fit=(BM) (WBM), "二部マッチング" {below, UniBlue}]{};
            \node[mysubset, AlertOrange, fill=AlertOrange!10, fit=(NF) (MCF), "ネットワークフロー" {above left, AlertOrange}]{};
            \node[mysubset, UniGreen, fill=UniGreen!10, fit=(MTR) (MI), "マトロイド" {below, UniGreen}]{};
        \end{scope}
    \end{tikzpicture}
\end{frame}

\begin{frame}[allowframebreaks]
    \frametitle{各回の内容（予定）} 
    \small
    \begin{enumerate}
        \item (4/23) イントロ+多面体的組合せ論（線形計画法の復習，整数多面体，完全単模行列）
        \item[] \textcolor{gray}{(4/30) 休み}
        \item (5/7) 二部マッチング①（Konig-Egervaryの定理，増加道アルゴリズム，ハンガリー法）
        \item (5/14) 二部マッチング②（最短路問題の復習，逐次最短路法と主双対法，最適性基準からの見方）
        \item (5/21) 最大流① （定式化，最大流最小カット定理，応用例）
        \item (5/28) 最大流②（残余ネットワーク，Ford--Fulkerson法，Edmonds--Karpのアルゴリズム）+ 最小費用流①（定式化，輸送問題，最大流との関係）
        \item[] \textcolor{gray}{(6/4) 休み}
        % \item[] \sout{(12/11) 最小カット（永持茨木のアルゴリズム，Kargerのアルゴリズム）}
        \item (6/11) 最小費用流②（逐次最短路法，容量スケーリング法）
        \item[] \textcolor{gray}{(6/18) 休み}
        \item[] \textcolor{gray}{(6/25) 休み}
        \item (7/2) マトロイド（定義と公理系，貪欲法，マトロイド多面体） 
        \item (7/9) 独立マッチング・マトロイド交差 （定義，応用例，Edmondsの最大最小定理，増加道アルゴリズム）
        \item (7/16) 劣モジュラ関数①（諸例，劣モジュラ基多面体，Lov\'asz拡張，劣モジュラ最小化）
        \item[] \textcolor{gray}{(7/23) 休み}
        \item (7/30) 劣モジュラ関数②（劣モジュラ最大化，近似アルゴリズム，貪欲法）
        \item (どこか) 精選トピック①
        \item (どこか) 精選トピック②
        \item (どこか) 精選トピック③
    \end{enumerate}

    \bigskip
    \begin{itemize}
        \item レポート課題を6月と期末に出す予定
        \item 精選トピックの日時は今後調整 
    \end{itemize}
\end{frame}